\section{存在的问题}

在本次凯尔纳目镜的设计与优化过程中，遇到了以下主要问题：

\subsection{（1）场曲校正与畸变的矛盾}

在优化过程中发现，场曲的减小往往伴随着畸变的增大。凯尔纳目镜三片结构的自由度有限（仅6个曲率半径和2个间距可调），很难同时将场曲和畸变都压缩到最低水平。为此，在优化时需要在两者之间取得平衡，通过调整Merit Function中各操作数的权重来控制优化方向。最终选择优先压缩畸变（将DIMX权重提高至10），以满足\(\le\)5\%的硬性指标，场曲则相对放宽。

从最终结果来看，切向场曲的残余值约为0.05 mm，弧矢场曲约为0.03 mm，虽未完全为零，但在目视光学系统中人眼具有一定的调焦能力，这一残余场曲在实际观测中不会造成明显影响。

\subsection{（2）出瞳距离与视场角的制约关系}

增大视场角有利于提升观测体验，但同时会使出瞳距离缩短，因为场镜需要更靠近焦平面来折转更大角度的光线。在本设计中，为了维持\(\ge\)10 mm的出瞳距离，视场角被限制在±20°。若需要进一步扩大视场角至±25°甚至±30°，需要引入更多的光学元件（如增加一片场镜，变为四片结构），这将超出本次设计的要求范围。

此外，出瞳距离的大小与场镜到双合镜的间距密切相关：增大间距可以增加出瞳距离，但同时会增大光束直径，对双合镜口径的要求更高，也会使系统总长增加。

\subsection{（3）像高未能完全达到指标}

在优化过程中，当尝试将全视场的像高提升时，全视场的畸变会显著增大，通过全局优化、锤形优化等多种手段均未能同时满足像高和畸变两项指标。分析原因，可能是初始结构的选取与全局最优结构存在一定偏差，导致优化陷入了局部最优解。

若后续有机会对设计进行改进，建议尝试从不同的初始结构出发（如参考更多专利结构），或使用Zemax的Global Search功能（而非仅依赖锤形优化）来寻找全局最优解。

\subsection{（4）玻璃材料选取的局限性}

本次设计中玻璃材料固定为BK7和F2，未对玻璃材料进行系统性优化。实际上，通过玻璃替换可以进一步改善色差和场曲的校正效果。由于发现该问题时已临近设计截止，未能对玻璃材料进行全面筛选。未来可考虑在Zemax中启用玻璃替换功能（Glass Substitution），从完整的玻璃库（如SCHOTT、CDGM）中搜索最优组合。

\subsection{（5）公差分析中灵敏度过高的问题}

公差分析结果表明，双合镜的倾斜公差（TETX/TETY）灵敏度较高，需要控制在±0.05°量级。然而，在大批量生产的流水线上，如手机光学产品这种对装配精度要求极高的场合，这一精度水平可能带来较高的装配成本。建议在设计方面进一步优化，通过调整公差敏感面的曲率半径来降低该面的倾斜灵敏度，或在结构设计上引入自定心机构来保证装配精度。
